\section{Equity-weighted doubly robust policy learning}\label{sec:method}

This section sets out the decision problem, the identifying assumptions and the estimator. The
estimator combines three ingredients that have so far been used separately: the climate-econometric
use of year-to-year weather variation as an exogenous shock \citep{dell2012,burke2015,hsiang2016};
doubly robust, cross-fitted scores for heterogeneous effects \citep{robinson1988,chernozhukov2018};
and welfare-maximising treatment choice \citep{manski2004,kitagawa2018,athey2021}, extended here to
welfare criteria that weight gains by the income of the people who receive them
\citep{atkinson1970,kitagawa2021}.

\subsection{Decision problem}

A planner, such as a development bank, a donor coalition or a regional body, can build climate-AI
capacity in a limited share $k$ of countries. Capacity means the research groups, data
infrastructure and skilled staff needed to develop and run AI-based adaptation services such as
early warning, agro-advisory and disaster-risk targeting. The planner chooses a rule
$\pi(Z_i)\in\{0,1\}$ based on observable country characteristics $Z_i$, and cares about how much
capacity reduces the income losses that hot years inflict on each part of the income distribution.

Let $i$ index countries, $d=1,\dots,10$ income deciles and $j$ observation periods (survey spells).
$Y_{idj}$ is the annualised growth of the mean income of decile $d$, $S_{ij}$ is a temperature shock,
$T_{ij}\in\{0,1\}$ indicates climate-AI capability, and $X_{ij}$ is a vector of country
characteristics. Suppressing indices, the outcome equation for decile $d$ is
\begin{equation}\label{eq:model}
Y_d = \beta_d(X,T)\,S + g_d(X,T) + \varepsilon_d, \qquad E[\varepsilon_d\mid S,X,T]=0 .
\end{equation}
Here $\beta_d(X,T)$ is the marginal effect of a $1^{\circ}$C shock on the income growth of decile
$d$, which is allowed to vary freely with $X$ and with capability, and $g_d$ is an unrestricted
function. The object of interest is the \emph{resilience effect} of capability,
\begin{equation}\label{eq:rho}
\rho_d(X) = \beta_d(X,1) - \beta_d(X,0),
\end{equation}
the change in decile $d$'s sensitivity to shocks when a country with characteristics $X$ has
climate-AI capability. When $S$ measures hot anomalies and $\beta_d<0$, $\rho_d>0$ means that
capability reduces the damage.

\subsection{Identification}

\begin{assumption}[Shock exogeneity]\label{as:shock}
Conditional on $(X,T)$, the de-trended temperature shock $S$ is independent of the potential
outcome errors $\varepsilon_d$.
\end{assumption}

\begin{assumption}[Selection on observables for sensitivity]\label{as:select}
Conditional on $X$, capability $T$ is independent of the pair of potential shock sensitivities
$\{\beta_d(X,0),\beta_d(X,1)\}$.
\end{assumption}

\begin{assumption}[Overlap]\label{as:overlap}
$e(X)=P(T=1\mid X)\in[\eta,1-\eta]$ for some $\eta>0$.
\end{assumption}

Assumption~\ref{as:shock} is the standard argument of the climate-econometrics literature: once a
country-specific trend is removed, year-to-year temperature variation is as good as random
\citep{dell2012,hsiang2016}. Assumption~\ref{as:select} is weaker than unconfoundedness of $T$
for income levels. Capability may well be correlated with growth through $g_d$; what we require is
that, among countries with the same income, inequality, connectivity, agricultural share, climate
and institutions, capability is not systematically held by countries whose incomes are more (or
less) shock-sensitive for unobserved reasons. We return to this assumption in the sensitivity
analysis.

\subsection{Estimator}

\paragraph{Step 1: cross-fitted nuisance functions.} Countries are split into $K=5$ folds. For
each fold, the following functions are estimated with gradient-boosted trees on the other folds:
$\hat s(X,T)=E[S\mid X,T]$; the conditional shock variance $\hat v(X,T)=E[(S-\hat s)^2\mid X,T]$;
the outcome mean $\hat\ell_d(X,T)=E[Y_d\mid X,T]$; and the capability propensity
$\hat e(X)=P(T=1\mid X)$, trimmed to $[0.05,0.95]$.

\paragraph{Step 2: slope scores.} Model~\eqref{eq:model} implies
$E[(Y_d-\ell_d)(S-s)\mid X,T]=\beta_d(X,T)\,v(X,T)$, so the Robinson-type ratio
$\psi_d=(Y_d-\hat\ell_d)(S-\hat s)/\hat v$ is an unbiased signal of $\beta_d(X,T)$. This signal is
insensitive to errors in $\hat\ell_d$ and $\hat s$ but not to errors in $\hat v$. We therefore
fit a preliminary slope function $\hat m_{d,t}(X)=E[\psi_d\mid X,T=t]$ and use the orthogonalised
score
\begin{equation}\label{eq:psistar}
\psi^{*}_d = \hat m_{d,T}(X) + \frac{(Y_d-\hat\ell_d)(S-\hat s) - \hat m_{d,T}(X)\,(S-\hat s)^2}{\hat v(X,T)} ,
\end{equation}
whose conditional mean is $\beta_d(X,T)$ and whose first-order sensitivity to every nuisance
function is zero.

\paragraph{Step 3: doubly robust resilience scores.} Augmented inverse-propensity weighting
of $\psi^{*}_d$ across capability states gives
\begin{equation}\label{eq:gamma}
\Gamma_d = \hat m_{d,1}(X) - \hat m_{d,0}(X)
 + \frac{T\,\{\psi^{*}_d-\hat m_{d,1}(X)\}}{\hat e(X)}
 - \frac{(1-T)\,\{\psi^{*}_d-\hat m_{d,0}(X)\}}{1-\hat e(X)} .
\end{equation}
$E[\Gamma_d\mid X]=\rho_d(X)$ if either the slope functions or the propensity score are correctly
specified, and $\Gamma_d$ is Neyman-orthogonal \citep{chernozhukov2018}. Averages of $\Gamma_d$
estimate decile-specific average resilience effects. Plotted against $d$, they form a
\emph{resilience incidence curve}, which shows whether capability protects the bottom of the
distribution more or less than the top.

\paragraph{Step 4: equity weighting.} Following \citet{atkinson1970}, the marginal social value of
a unit of income growth for decile $d$ of country $i$ is
$\omega_{id}(\epsilon)=(y_{id}/\bar y)^{-\epsilon}$, where $y_{id}$ is the decile's mean
consumption or income in purchasing-power-parity dollars, $\bar y$ the world population-weighted
mean, and $\epsilon\geq0$ the planner's aversion to inequality. The weights are normalised to a
population-weighted mean of one, so $\epsilon=0$ reproduces the utilitarian criterion. The
equity-weighted benefit of building capacity in country $i$ is
\begin{equation}\label{eq:benefit}
B_i(\epsilon) = \bar s_i \sum_{d=1}^{10} p_d\,\omega_{id}(\epsilon)\,\bar\Gamma_{id},
\end{equation}
where $\bar\Gamma_{id}$ averages the scores over country $i$'s spells, $p_d=0.1$ is each decile's
population share, and $\bar s_i$ is the country's typical shock size (the standard deviation of its
de-trended temperature). Because $B_i(\epsilon)$ is a fixed linear combination of orthogonal scores,
it is itself an orthogonal score for the equity-weighted benefit
$b_i(\epsilon)=\bar s_i\sum_d p_d\,\omega_{id}(\epsilon)\,\rho_d(X_i)$.

\paragraph{Step 5: policy learning.} The planner's rule is chosen from $\Pi$, the class of depth-2
decision trees on a short list of interpretable characteristics $Z$, by solving
\begin{equation}\label{eq:policy}
\hat\pi_{\epsilon} = \arg\max_{\pi\in\Pi}\; \frac{1}{N}\sum_{i=1}^{N} w_i\,\pi(Z_i)\,B_i(\epsilon)
\quad\text{s.t.}\quad \frac{1}{N}\sum_{i=1}^{N}\pi(Z_i)\le k ,
\end{equation}
with country population weights $w_i$. Depth-2 trees keep the rule readable as a short list of
conditions, which is what makes it usable as a policy instrument, and the problem can be solved
exactly by enumerating splits at the deciles of each feature together with every budget-feasible
set of leaves.

\begin{proposition}[Regret]\label{prop:regret}
Suppose Assumptions~\ref{as:shock}--\ref{as:overlap} hold, the welfare weights are bounded,
$0<\omega_{id}(\epsilon)\le\bar\omega_\epsilon$, shock sizes are bounded, $\bar s_i\le\bar s$, the
nuisance estimators are consistent with
$\|\hat e-e\|_2\cdot\|\hat m_{d,t}-\beta_{d}(\cdot,t)\|_2=o_p(N^{-1/2})$ and analogous product-rate
conditions for $(\hat s,\hat v,\hat\ell_d)$, and $\Pi$ has finite VC dimension. Then the regret of
$\hat\pi_\epsilon$ relative to the best rule in $\Pi$ under the $\epsilon$-weighted criterion
satisfies
\[
\sup_{\pi\in\Pi}W_\epsilon(\pi)-W_\epsilon(\hat\pi_\epsilon)
=O_p\!\left(\bar s\,\bar\omega_\epsilon\,\sqrt{\mathrm{VC}(\Pi)/N}\right),
\qquad W_\epsilon(\pi)=E[w\,\pi(Z)\,b(\epsilon)] .
\]
\end{proposition}

\noindent\emph{Proof sketch.} $B_i(\epsilon)$ is a bounded linear combination of the orthogonal
scores $\Gamma_{id}$ with coefficients $\bar s_i p_d\,\omega_{id}(\epsilon)$ that do not depend on
the nuisance functions. Hence it inherits Neyman orthogonality and the product-rate condition makes
the error from estimating the nuisances $o_p(N^{-1/2})$ uniformly over $\Pi$. The remaining term
is the empirical process of an oracle score whose envelope is scaled by $\bar s\,\bar\omega_\epsilon$,
and the bound follows from Theorem~1 of \citet{athey2021} applied to that scaled score. The budget
constraint restricts $\Pi$ to a subset and does not increase its complexity. \hfill$\square$

The proposition makes one trade-off explicit. Greater inequality aversion raises
$\bar\omega_\epsilon$, because the poorest deciles in the poorest countries receive large weights,
so the statistical price of an equity-minded rule is a larger regret bound. With $N$ in the low
hundreds this price is not negligible, and the Monte Carlo evidence in Section~\ref{sec:mc}
quantifies it.

\subsection{Evaluation and inference}

Policies are evaluated out of sample. The tree is learned on four country folds and its value is
computed from the held-out scores of the fifth, rotating across folds. We compare the
equity-weighted rule $\hat\pi_\epsilon$ with (i) the utilitarian rule $\hat\pi_0$ evaluated under
the same $\epsilon$; (ii) a plug-in rule that ranks countries by $\hat m_{d,1}-\hat m_{d,0}$ without
the doubly robust correction; (iii) an ``AI-to-the-capable'' rule that allocates capacity to the
countries with the most AI research; (iv) a ``poorest-first'' rule; and (v) random allocation.
Confidence intervals for policy values and their differences come from a country-cluster bootstrap
of the held-out values. Two summary quantities connect the method to the special issue's
question. The \emph{equity gain}, $W_\epsilon(\hat\pi_\epsilon)-W_\epsilon(\hat\pi_0)$, measures how
much welfare an equity-minded planner loses by using the utilitarian rule. The \emph{equity price},
$W_0(\hat\pi_0)-W_0(\hat\pi_\epsilon)$, measures the aggregate income resilience given up by
prioritising the poor.

Alongside the policy results, we report the average resilience effect and its best linear
projection on $Z$, and sorted group average treatment effects with classification analysis
\citep[GATES and CLAN;][]{chernozhukov2025}. These describe which characteristics separate the
countries in which capability matters most from those in which it does not.
